\section{Earlier synthetic and byte-language-model illustrations}
\label{app:earlier-illustrations}

This appendix preserves the earlier studies, separately from the pretrained-model study in Section~\ref{sec:experiments}. The theory is the main evidence for the global statements. These earlier experiments
check the response formula, illustrate finite reconstruction, and expose
precision and conditioning limits. They use the exact-product design of
Corollary~\ref{cor:response_tomography} or rank-checked natural/Gaussian
probes, \emph{not} the repeated-complement design of
Corollary~\ref{cor:finite-noisy-main}. Thus they do not empirically evaluate
its whole-feasible-set diameter bound. All numerical tables come from
validated result records; protocols, negative results, source hashes, and
replay limitations are retained in
Appendices~\ref{sec:benchmark}--\ref{app:trusted-experiments}.

\paragraph{Independent full-dimensional response measurements.}
Three soft-router byte-LM checkpoints have $L=12$, $K=4$, and $D=787{,}712$.
We measure responses by reverse-mode differentiation of a linear effective
loss followed by a forward-mode JVP of the packed factor-to-product map.
The observation code does not call Equation~\ref{eq:response_blocks}.
For each native, permutation, and fixed non-permutation chart, four designed
probes reconstruct the response and two unseen Gaussian probes test its
predictions. Table~\ref{tab:autograd-tomography} shows numerical agreement
with the formula, successful held-out prediction, and separation of the
fixed alternative. This finite check does not prove separation of every
possible alternative; that quantifier is supplied by the theorem.

% Generated by experiments/summarize_autograd_tomography.py.
\begin{table}[H]
\centering\small
\caption{Full-dimensional autodiff response audit. Errors are relative Frobenius norms. Formula and held-out columns take the maximum over native, permutation, and non-permutation charts; each chart uses four designed and two unseen probes.}
\label{tab:autograd-tomography}
\begin{tabular}{@{}crrrr@{}}
\toprule
Seed & Formula error & Held-out error & Permutation gap & Non-perm. gap \\
\midrule
0 & 3.52e-16 & 4.68e-14 & 9.73e-17 & 0.9734 \\
1 & 3.37e-16 & 5.07e-14 & 1.05e-16 & 0.9731 \\
2 & 3.71e-16 & 4.12e-14 & 1.11e-16 & 0.9735 \\
\bottomrule
\end{tabular}
\end{table}


\paragraph{Natural directions and finite precision.}
Actual byte-LM cross-entropy gradients with respect to independent effective
layer weights supply 16 directions per checkpoint. Controlled autodiff
responses are measured at a fixed checkpoint; the first $q\in\{1,4,8,12\}$
directions fit the structural estimator and the last four test prediction.
With the exact product row space, four directions meet the sufficient
reconstruction ranks at all six checkpoints in
Table~\ref{tab:natural-gradients}. Factor errors use permutation alignment
only for evaluation. One direction has insufficient local rank and gives
errors $0.258$--$1.34$ at the three longer-trained checkpoints.

% Keep the generated second table after the first without changing its source.
\begingroup
\let\originaltable\table
\renewcommand{\table}[1][]{\originaltable[H]}
\begin{table}[t]
\centering\small
\caption{Four natural task-gradient probes with known response structure. Original and new checkpoints use the separately registered 1,500- and 6,000-step configurations. Recovery error is the combined relative factor error after evaluation-only permutation alignment. These are FP64 fixed-checkpoint measurements, not observed training trajectories.}
\label{tab:natural-gradients}
\begin{tabular}{llrr}
\toprule
Configuration & Seed & Max. held-out response error & Factor error \\
\midrule
Original & 0 & $9.57\times10^{-14}$ & $3.02\times10^{-8}$ \\
Original & 1 & $8.22\times10^{-14}$ & $5.37\times10^{-9}$ \\
Original & 2 & $8.63\times10^{-14}$ & $1.36\times10^{-8}$ \\
New & 10 & $9.39\times10^{-14}$ & $2.31\times10^{-8}$ \\
New & 11 & $6.45\times10^{-14}$ & $1.35\times10^{-8}$ \\
New & 12 & $9.95\times10^{-14}$ & $1.48\times10^{-8}$ \\
\bottomrule
\end{tabular}
\end{table}

\endgroup

Finite Euclidean steps reveal an arithmetic limit. At step size $10^{-5}$,
the longer-trained checkpoints have FP64 response errors
$3.3\times10^{-8}$--$5.2\times10^{-8}$, but FP32 subtraction errors
$0.0208$--$0.0260$ and BF16 errors about one. These tests cast factors and
arithmetic to the stated precision; they do not test mixed-precision
training with FP32 master weights. The step-size sweep is in
Appendix~\ref{app:enhancement-experiments}.

\paragraph{Noisy recovery does not supply an automatic trust certificate.}
A separate frozen development/calibration/evaluation protocol compares
spectral recovery, simplex projection with a basis refit, and three-start
constrained joint least squares. Both the product and measured responses
receive relative noise up to $10^{-2}$; the estimated product subspace is
recomputed. On 1,920 held-out synthetic variants across 40 seeds, joint
fitting reduces the median factor error from $3.34\times10^{-3}$ for
projected spectral recovery to $1.01\times10^{-4}$, and errors above $0.10$
from 566 to 227. All 48 variants at four fresh byte-LM checkpoints have
joint-fit errors below $0.0107$.

Failures remain substantial: the largest joint-fit error is $10.54$, and a
combined diagnostic falsely accepts an error of $0.82$. Its accepted-bad
fraction is $92/1782=5.16\%$, with seed-cluster bootstrap interval
$[4.50\%,5.79\%]$; residual-only rejection yields $68/1757=3.87\%$.
The combined score therefore does not establish superiority, and its point
estimate exceeds the $5\%$ calibration target. The candidate-local checks
also do not exclude distant solutions or establish truth membership in the
fitted subspace. Appendix~\ref{app:trusted-experiments} retains all methods,
tail errors, failed solver statuses, and independent CPU replay differences.
None of these diagnostics is the global certificate above.

\paragraph{Auxiliary structural decisions.}
Identifying coordinates under an optimizer does not make actions derived
from them product invariant. In a documented reimplementation of ASLoRA's
first RoBERTa/MRPC merge rule, a fixed grid changes the primary action in
$11/30$ seed--gauge instances. Only the selected layer indices are carried
back to a common canonical model; validation-loss differences have both
signs. A separate controlled byte-LM partition audit finds changes at all
three longer-trained checkpoints, but matched recovery training greatly
reduces the performance gaps. These are bounded decision witnesses, not
applications of the softmax theorem to LoRA, universal harm claims, or
full-training regret evaluations. Detailed results, failed signed-local
searches, and the earlier initialization-dominated experiment remain in
Appendices~\ref{app:detailed-results} and~\ref{app:enhancement-experiments}.
